\subsection{Theory as guide towards generalizability}

A core function of economics is to construct models that capture causal and generalizable relationships that remain stable across environments.
For example, the idea that humans make reference-dependent utility choices is not specific to one particular economic environment, but has been shown to broadly apply to decision-making under uncertainty \citep{prospectkahneman1979}.
Conveniently, these are exactly the types of generalizable relationships that we would expect to better predict human behavior in new settings when supplied to an LLM.
Our approach is to narrow the search space of possible \persona{s} by grounding candidate \persona{s} in such theories.
By doing so, we might increase the likelihood of identifying \persona{s} that we have prior reason to believe reflect genuine, stable behavioral patterns.
Without doing so, we risk dramatically underfitting and failing even to accurately predict the training data.

Economic theories---and theories from other social sciences that embrace methodological individualism \citep{friedman1953methodology}---are well-suited to be cast into \persona{s}, since good theory rests on the choices and actions of individuals.
Specifically, we consider a prompt to be ``theoretically grounded'' when it instructs the LLM to make predictions about how a human would respond based on some underlying theory or model of human behavior.
For example, a utility function that trades off one's own payoff against inequality might become \emph{``You value your own earnings but dislike outcomes where you earn far more (or less) than others.''}
A highly capable and tool-using LLM could even be given that utility function and parameters directly (e.g., $u(x_{self}, x_{other})= x_{self} - | x_{self} - x_{other}|$).
A prospect-theoretic model reasonably maps to \emph{``You are risk averse in gains, risk seeking in losses, and probability weight the very likely and very unlikely outcomes''}---and of course these could be parameterized as well.

This definition is necessarily broad because what constitutes a theory is also broad.
One might consider good theories to be those that are parsimonious, predictive, and interpretable.
We view these as effective guides for constructing \persona{s} that are grounded in theory.
But just as there is no universal rule as to what makes a theory good, there is not a singular all-encompassing playbook for candidate \persona{s}.
Here, the researcher must leverage their expertise appropriately to determine which theories are most relevant to the domain where predictions are desired.
This should be relatively straightforward in practice.
One can even simply ask an LLM to construct a candidate set of prompts based on a given paper and then adjust the prompts accordingly.
For example, the linked notebook---which took 15 minutes to code---contains an example where we take PDFs of 5 well-known papers in the behavioral economics literature and generate 10 agents from each.\footnote{\url{https://expectedparrot.com/content/32751dae-7a69-430d-9f1a-0c3f50ccef5b}}

It is worth highlighting just how flexible such agents are as predictive models.
They can make predictions in response to \emph{any} setting described in natural language.
This is usually impossible for traditional machine learning and mathematical economic models, which operate within a fixed parameter space. 
One cannot introduce additional covariates to a regression model once it is trained. 
Similarly, a prospect-theoretic model for evaluating risky gambles cannot instantaneously incorporate other contextually relevant factors---the recent performance of financial markets, the demographic composition of the experimental sample, or even something as mundane as whether participants are making decisions before or after lunch. 
\aisub{s}, by contrast, can interpolate over such details through natural language instructions. 
Whether these agents actually leverage this flexibility effectively is, of course, an empirical question, but one that is easily testable.

Once a set of \persona{s} is defined, we can optimize them against an appropriate sample of relevant human decisions.
This may involve searching for an optimal mixture of \persona{s}, or adding continuous parameters directly in a single prompt and adjusting them until the simulated distribution aligns with the training data.
From a machine learning perspective, the candidate set of \persona{s} defines the functional form of our hypothesis class, with theory guiding this specification to effectively navigate the bias-variance tradeoff.\footnote{One could imagine an ``oracle'' \persona{} akin to a perfect program, describing every preference, heuristic, and belief update rule---allowing the LLM to accurately produce responses in all contexts for a person.
In effect, our approach is to identify portions of this oracle that are most relevant to the given training and testing data.
Increasing context windows suggest that such highly generalizable agents may be possible in the not-so-distant future.}
The set is then pruned (or tuned, depending on the method) to best fit the training data.

As an example, consider the extensive theoretical and empirical economic literature studying how people reason strategically \citep{STAHL1994,STAHL1995,1120Arad2012, Camerer2004Cognitive}.
Level-$k$ models, for instance, predict that individuals best respond based on their beliefs about others' levels of reasoning. 
In the $\frac{2}{3}$ guessing game, players aim to pick $\frac{2}{3}$ of the average number chosen by the group \citep{LkNagel1995,keynes1936general}. 
If we had human data from such a game, we could construct a series of \persona{s} that specify different levels of strategic thinking (e.g., \emph{``You are a level-0 reasoner,'' ``You are a level-1 reasoner,''} etc.).
We then identify the mixture that best fits the training distribution and test whether it generalizes to other variants of the game (e.g., different values than $\frac{2}{3}$).
Then, the optimized \persona{s} can be used to predict other related distributions in new settings.

Of course, constructing and validating these \persona{s} is not a perfectly specified problem.
There are various ways to translate a theory into natural language, and multiple theories may apply to a given setting.
The boundaries of a theory and the settings it plausibly governs are not always well-defined.
Nor can we guarantee that the data-generating processes of the training and testing sets are either sufficiently similar or sufficiently distinct to yield reliable validation.
Yet, as we will show empirically, these challenges can be overcome in practice, at least in some domains.

Ultimately, applying a \persona{}---or a set of \persona{s}---to a new environment requires an unavoidable inductive leap.
What we can do is try to make the leap explicit and interpretable.  
LLMs are extremely good at following explicit, well-defined instructions.
Because each \persona{} is a set of these flexible natural language instructions, researchers can both evaluate performance and reasonably assess its relevance for a new setting.\footnote{Such flexibility also makes \aisub{s} less susceptible to the Lucas critique \citep{Lucas1976Critique}---if at all.
Unlike traditional agent-based models or mathematical theories, which have hard-coded predetermined processes to generate predictions, LLMs can, as we will show, interpret new policies in context, reason through their implications, and adaptively formulate responses.
}
This mirrors how economic theory is used more broadly with real humans: it is tested against data, observed where it succeeds or fails, and cautiously extrapolated to settings just beyond current evidence.
When a new tariff is proposed, for example, the theory motivating the tariff is not known to be universally ``correct'' in a complex dynamic world ex ante.
It is a disciplined guess, shaped by assumptions and previous evidence, which the economist then uses to make predictions about the world.
Theory-guided \persona{s}, validated across environments, bring the same kind of disciplined reasoning to simulated \aisub{} predictions of human behavior.
